% Pursuing Stacks (cotes 134-2 to 134-8): archive cote 134-7, site page 87 (PDF page 88), typescript page 529.
% typescript. Draft compared: ps6.tex:4373-4449.
% AI-transcribed by Claude Opus 5.5 (claude-opus-5-5, adaptive thinking, effort high), Message Batches API, batch ps4 (msgbatch_012Aoc1NuskATPsyBeKSYK7J), 2026-11-12; not yet checked by eye.
\scanpage{134-7/87}
\mspage{529}

\inserted{functor,} though is lost on our way - I am affraid I am confusing the kind of dua-%
lity I am after here, with a rather different type of duality, kin to
Poincaré duality, and \inserted{discovered} \xout{\illegible{1}} I believe by J.H.C. Whitehead, in the context
of spaces having the homotopy type of a finite complex. For such a
complex, the idea \inserted{(as far as I remember)} is to embed X into a large-dimensional sphere, and
to take the complement $X\hand{^{\vee}}$ of an open tubular neighbourhood of X. Up to
suspension, the homotopy type of X does not depend on the choices made,
and \inserted{(shifting back by $n$)} we thus get a canonical object in the suspension category, depending
contravariantly on X. The functor $X\hand{\mapsto}X\hand{^{\vee}}$ (if I got it right) is an auto-%
duality of the relevant full subcategory of the suspension category,
and the (reduced) homology functor $L\overset{\hand{\sim}}{H}.$ \xout{\unsure{commutes to}} maps this subcategory
into perfect complexes, and \unsure{Commutes|commutes} to the autodualities. The whole
story seems tailored towards a study of duality relations for the functor
$L\overset{\hand{\sim}}{H}.$ exclusively - without any reference to the homotopy invariants $\hand{\pi}_{i}$.
I do'nt know if one can devise a similar story for $L\hand{\pi}.$ , by stabilizing
with respect to the loop space functor (or is this just nonsense ?). At
any rate, there does'nt seem any autoduality in view, exchanging the
homology and homotopy invariants ...

E) \uline{Conservativity properties:} The functors $L\overset{\hand{\sim}}{H}.$ and $L\hand{\pi}.$ are both
conservative, i.e. a map in $\mathrm{Hot}_{\bullet}(k)$ \inserted[typed]{which} \xout{mapped} by either of these functors
becomes an isomorphism, is an isomorphism. (Here of course the 0-connec%
-tedness assumption for the \xout{\unsure{objects we are}} homotopy types we are working
with is essential, as far as the functor $L\hand{\pi}.$ is concerned at any rate.)

F) \uline{Base change properties}: All four functors, and the adjunction
and inversion maps (7) and (9), are compatible with ring extension
$k\hand{\to} k'$ , it being understood that such a ring homomorphism defines
functors
\[ (12)\qquad \mathrm{Hotab}_{\bullet}(k)\hand{\longrightarrow}\mathrm{Hotab}_{\bullet}(k') \ , \quad \mathrm{Hot}_{o}(k)\hand{\longrightarrow}\mathrm{Hot}_{\bullet}(k') \]
compatible with the fibration and cofibration structures.

In opposite direction, there should be too \inserted{a} ``restriction'' functor\xout{s}
(less important though I feel), I did'nt try to find out what \xout{their} \inserted{it's}
formal properties should be with respect to the formalism developped her\unsure{e}

G) \uline{Comparison with ordinary homotopy types }: We got functors
\[ (13)\qquad \mathrm{Hotab}_{\bullet}(k)\hand{\longrightarrow} \mathrm{Hotab} = D.(Ab) \ , \quad \mathrm{Hot}_{o}(k)\hand{\longrightarrow} \mathrm{Hot}_{\bullet} \quad , \]
(the first being interpreted as ``forget k'' functor, the second \deleted{\illegible{1}} \inserted{as a}
sections functor on semisimplicial schematic models). This pair of

\margin{87}
